% Created 2026-05-11 Mon 12:29
% Intended LaTeX compiler: pdflatex
\documentclass[11pt]{article}

\usepackage[utf8]{inputenc}
\usepackage[T1]{fontenc}
\usepackage{graphicx}
\usepackage{longtable}
\usepackage{wrapfig}
\usepackage{rotating}
\usepackage[normalem]{ulem}
\usepackage{amsmath}
\usepackage{amssymb}
\usepackage{capt-of}
\usepackage{hyperref}
\author{flavio}
\date{\today}
\title{Introduzione corso}
\hypersetup{
 pdfauthor={flavio},
 pdftitle={Introduzione corso},
 pdfkeywords={},
 pdfsubject={},
 pdfcreator={},
 pdflang={English}}
\begin{document}

\maketitle
\tableofcontents

\section{Programmazione orientata agli oggetti}
\label{sec:org7a8805e}
La progrmmazione orientata agli oggetti fornisce nuovi elementi per rappresentare elementi nello spazio del problema:
\begin{itemize}
\item Gli elementi sono chiamati \textbf{oggetti}
\begin{itemize}
\item Pensiamo a modellare gli oggetti, non il programma (che vuol dire?)
\end{itemize}
\end{itemize}
\subsection{Tutto è un oggetto:}
\label{sec:org7a456f2}
\begin{itemize}
\item Ha uno stato con cui altri oggetti o sé stesso interagiscono.
\begin{itemize}
\item Ha delle operazioni con che puoi chiedergli di eseguire
Per effettuare una richiesta ad un oggetto si manda un messaggio a quell'oggetto. Un messaggio è una \textbf{chiamata} a una \textbf{funzione} (metodo) che appartiene a un particolare oggetto.
\end{itemize}
\end{itemize}
\subsection{Incapsulamento e information hiding}
\label{sec:orgb80341d}
Un programma può nascondere la sua complessità mediante la semplicità degli oggetti
\subsection{Ogni oggetto possiede un tipo detto \textbf{classe} ed è istanza di una classe. La classe viene identificata dai messaggi (metodi).}
\label{sec:org00b96e2}
\subsection{Ereditarietà: vogliamo evitare di ricreare nuove classi di oggetti quando esse hanno funzionalità simili.}
\label{sec:orga168ba2}
\subsection{Polimorfismo:Basta sapere la classe base, senza dover conoscere la classe specifica.}
\label{sec:orga114803}
\section{Il linguaggio che useremo: Java}
\label{sec:orgfd3026e}
Java un linguaggio creato all'alba del nuovo millennio: orientato agli oggetti.
\subsection{Caratteristiche principali}
\label{sec:orgc10214c}
\begin{itemize}
\item Portabile: WORA(\textbf{write once, run anywhere})
Al contrario dl linguaggi come il  C o C++ non viene compilato per una piattaforma, ma viene compilato in \textbf{bytecode}, per una macchina virtuale la \textbf{Java virtual machine}. Il bytecode è tradotto al volo in istruzioni per la piattaforma.
\item Sicurezza: La maggior parte degli errori che possono avvenire a run-time vengono eliminate a tempo di compilazione. Dove non è possibile sono gestite a tempo di esecuzione.
\end{itemize}
\end{document}
